\subsection{Knowledge extraction evaluation set}\label{subsec:knowledge-extraction-eval-set}

The calibration of the knowledge-extraction pipeline thresholds, matcher settings, and the agreement-based cascade configuration is done, on the calibration dataset of Subsection~\ref{subsec:knowledge-extraction-calibration-set}. To evaluate the performance of the knowledge-extraction pipeline and measure strict-F1 scores, a held-out corpus of 15 papers is used. Once the calibration is complete, no thresholds or hyperparameters are changed on the evaluation papers. 

The evaluation is done on a set of 15 papers, randomly and manually selected from the ingested histopathology corpus of Subsection~\ref{subsec:corpus-source-and-scope}. This evaluation set is fully disjoint from the calibration set. Unlike the calibration cluster, the papers are chosen randomly, not based on relatedness. This is intentional: the evaluation measures per-chunk extraction quality, which does not require cross-paper signals; a random sample is expected to be more generalizable to unseen papers, compared to a subset of related papers. 
This held-out set is the basis of the generalization test of RQ5 (Subsection~\ref{subsec:eval-heldout-generalization}).

The evaluation set is processed in a way similar to the calibration set: each paper is segmented into MAP-stage chunks, and every usable paragraph becomes one case: all paragraphs that are at least 40 words long, and are outside of references, acknowledgments, and similar non-content sections. For each source case, silver findings are generated by prompting Opus 4.7, yielding 285 source cases, and from these cases, 1\,657 silver findings.
